% Project Journey Portfolio — group submission (one PDF).
% Build: latexmk -pdf project-journey-portfolio.tex
\documentclass[11pt,a4paper]{article}
% Shared style for the StudyOS Thesis-Finder course submissions.
% Included by every submission document via % Shared style for the StudyOS Thesis-Finder course submissions.
% Included by every submission document via % Shared style for the StudyOS Thesis-Finder course submissions.
% Included by every submission document via \input{preamble}.
% Plain black-and-white house style: no colour, prints identically in mono.

\usepackage[T1]{fontenc}
\usepackage[utf8]{inputenc}
\usepackage{lmodern}
\usepackage[english]{babel}
\usepackage{microtype}
\usepackage[a4paper,margin=2.4cm]{geometry}
\usepackage{xcolor}
\usepackage{titlesec}
\usepackage{enumitem}
\usepackage{booktabs}
\usepackage{tabularx}
\usepackage{tcolorbox}
\tcbuselibrary{breakable,skins}
\usepackage{pifont}
\usepackage{xurl}
\usepackage[hidelinks]{hyperref}

\urlstyle{tt}

% --- Palette (plain black and white) ----------------------------------------
% Everything is black on white; the only greys are for a hairline rule and the
% callout title bar, so the documents print identically in mono.
\definecolor{sosAccent}{gray}{0}
\definecolor{sosDark}{gray}{0}
\definecolor{sosOpen}{gray}{0}
\definecolor{sosOpenBg}{gray}{1}
\definecolor{sosRule}{gray}{0.6}

% --- Headings ---------------------------------------------------------------
\titleformat{\section}{\normalfont\large\bfseries}{\thesection}{0.6em}{}
\titleformat{\subsection}{\normalfont\normalsize\bfseries\itshape}{\thesubsection}{0.6em}{}
\titlespacing*{\section}{0pt}{1.1\baselineskip}{0.35\baselineskip}
\titlespacing*{\subsection}{0pt}{0.7\baselineskip}{0.2\baselineskip}

\setlist[itemize]{leftmargin=1.2em,itemsep=0.15em,topsep=0.25em,parsep=0pt}
\setlist[enumerate]{leftmargin=1.4em,itemsep=0.15em,topsep=0.25em,parsep=0pt}

% --- Semantic helpers -------------------------------------------------------
% Inline reference to a file or path inside the project repository.
\newcommand{\repofile}[1]{\path{#1}}

% One "the assignment asks for ..." lead-in line.
\newcommand{\asks}[1]{\noindent\textbf{Assignment asks for:} #1\par\vspace{0.3em}}

% Open question the team still has to answer (was <mark>highlighted</mark> in the
% markdown draft).
\newenvironment{openbox}[1][Open --- needs the team]{%
  \begin{tcolorbox}[
    breakable,
    colback=white,
    colframe=black,
    colbacktitle=black!8,
    coltitle=black,
    boxrule=0.4pt,
    titlerule=0pt,
    left=0.6em,right=0.6em,top=0.5em,bottom=0.5em,
    title={#1},
    fonttitle=\bfseries\small,
    sharp corners,
  ]\small
}{%
  \end{tcolorbox}
}

\NewDocumentEnvironment{infobox}
  { O{Open --- needs the team} O{} }
{%
  \begin{tcolorbox}[
    breakable,
    colback=white,
    colframe=black,
    colbacktitle=black!8,
    coltitle=black,
    boxrule=0.4pt,
    titlerule=0pt,
    left=0.6em,
    right=0.6em,
    top=0.5em,
    bottom=0.5em,
    title={#1},
    fonttitle=\bfseries\small,
    sharp corners,
    #2
  ]\small
}{%
  \end{tcolorbox}
}


% Status markers for the submission checklist.
\newcommand{\stdone}{\ding{51}}
\newcommand{\stpartial}{\ding{108}}
\newcommand{\stgap}{\ding{55}}

% Title block.
\newcommand{\sostitle}[3]{%
  \begin{center}
    {\LARGE\bfseries #1}\\[0.35em]
    {\large\itshape #2}\\[0.5em]
    {\small #3}
  \end{center}
  \vspace{0.4em}
  {\color{sosRule}\hrule height 0.4pt}
  \vspace{0.9em}
}
.
% Plain black-and-white house style: no colour, prints identically in mono.

\usepackage[T1]{fontenc}
\usepackage[utf8]{inputenc}
\usepackage{lmodern}
\usepackage[english]{babel}
\usepackage{microtype}
\usepackage[a4paper,margin=2.4cm]{geometry}
\usepackage{xcolor}
\usepackage{titlesec}
\usepackage{enumitem}
\usepackage{booktabs}
\usepackage{tabularx}
\usepackage{tcolorbox}
\tcbuselibrary{breakable,skins}
\usepackage{pifont}
\usepackage{xurl}
\usepackage[hidelinks]{hyperref}

\urlstyle{tt}

% --- Palette (plain black and white) ----------------------------------------
% Everything is black on white; the only greys are for a hairline rule and the
% callout title bar, so the documents print identically in mono.
\definecolor{sosAccent}{gray}{0}
\definecolor{sosDark}{gray}{0}
\definecolor{sosOpen}{gray}{0}
\definecolor{sosOpenBg}{gray}{1}
\definecolor{sosRule}{gray}{0.6}

% --- Headings ---------------------------------------------------------------
\titleformat{\section}{\normalfont\large\bfseries}{\thesection}{0.6em}{}
\titleformat{\subsection}{\normalfont\normalsize\bfseries\itshape}{\thesubsection}{0.6em}{}
\titlespacing*{\section}{0pt}{1.1\baselineskip}{0.35\baselineskip}
\titlespacing*{\subsection}{0pt}{0.7\baselineskip}{0.2\baselineskip}

\setlist[itemize]{leftmargin=1.2em,itemsep=0.15em,topsep=0.25em,parsep=0pt}
\setlist[enumerate]{leftmargin=1.4em,itemsep=0.15em,topsep=0.25em,parsep=0pt}

% --- Semantic helpers -------------------------------------------------------
% Inline reference to a file or path inside the project repository.
\newcommand{\repofile}[1]{\path{#1}}

% One "the assignment asks for ..." lead-in line.
\newcommand{\asks}[1]{\noindent\textbf{Assignment asks for:} #1\par\vspace{0.3em}}

% Open question the team still has to answer (was <mark>highlighted</mark> in the
% markdown draft).
\newenvironment{openbox}[1][Open --- needs the team]{%
  \begin{tcolorbox}[
    breakable,
    colback=white,
    colframe=black,
    colbacktitle=black!8,
    coltitle=black,
    boxrule=0.4pt,
    titlerule=0pt,
    left=0.6em,right=0.6em,top=0.5em,bottom=0.5em,
    title={#1},
    fonttitle=\bfseries\small,
    sharp corners,
  ]\small
}{%
  \end{tcolorbox}
}

\NewDocumentEnvironment{infobox}
  { O{Open --- needs the team} O{} }
{%
  \begin{tcolorbox}[
    breakable,
    colback=white,
    colframe=black,
    colbacktitle=black!8,
    coltitle=black,
    boxrule=0.4pt,
    titlerule=0pt,
    left=0.6em,
    right=0.6em,
    top=0.5em,
    bottom=0.5em,
    title={#1},
    fonttitle=\bfseries\small,
    sharp corners,
    #2
  ]\small
}{%
  \end{tcolorbox}
}


% Status markers for the submission checklist.
\newcommand{\stdone}{\ding{51}}
\newcommand{\stpartial}{\ding{108}}
\newcommand{\stgap}{\ding{55}}

% Title block.
\newcommand{\sostitle}[3]{%
  \begin{center}
    {\LARGE\bfseries #1}\\[0.35em]
    {\large\itshape #2}\\[0.5em]
    {\small #3}
  \end{center}
  \vspace{0.4em}
  {\color{sosRule}\hrule height 0.4pt}
  \vspace{0.9em}
}
.
% Plain black-and-white house style: no colour, prints identically in mono.

\usepackage[T1]{fontenc}
\usepackage[utf8]{inputenc}
\usepackage{lmodern}
\usepackage[english]{babel}
\usepackage{microtype}
\usepackage[a4paper,margin=2.4cm]{geometry}
\usepackage{xcolor}
\usepackage{titlesec}
\usepackage{enumitem}
\usepackage{booktabs}
\usepackage{tabularx}
\usepackage{tcolorbox}
\tcbuselibrary{breakable,skins}
\usepackage{pifont}
\usepackage{xurl}
\usepackage[hidelinks]{hyperref}

\urlstyle{tt}

% --- Palette (plain black and white) ----------------------------------------
% Everything is black on white; the only greys are for a hairline rule and the
% callout title bar, so the documents print identically in mono.
\definecolor{sosAccent}{gray}{0}
\definecolor{sosDark}{gray}{0}
\definecolor{sosOpen}{gray}{0}
\definecolor{sosOpenBg}{gray}{1}
\definecolor{sosRule}{gray}{0.6}

% --- Headings ---------------------------------------------------------------
\titleformat{\section}{\normalfont\large\bfseries}{\thesection}{0.6em}{}
\titleformat{\subsection}{\normalfont\normalsize\bfseries\itshape}{\thesubsection}{0.6em}{}
\titlespacing*{\section}{0pt}{1.1\baselineskip}{0.35\baselineskip}
\titlespacing*{\subsection}{0pt}{0.7\baselineskip}{0.2\baselineskip}

\setlist[itemize]{leftmargin=1.2em,itemsep=0.15em,topsep=0.25em,parsep=0pt}
\setlist[enumerate]{leftmargin=1.4em,itemsep=0.15em,topsep=0.25em,parsep=0pt}

% --- Semantic helpers -------------------------------------------------------
% Inline reference to a file or path inside the project repository.
\newcommand{\repofile}[1]{\path{#1}}

% One "the assignment asks for ..." lead-in line.
\newcommand{\asks}[1]{\noindent\textbf{Assignment asks for:} #1\par\vspace{0.3em}}

% Open question the team still has to answer (was <mark>highlighted</mark> in the
% markdown draft).
\newenvironment{openbox}[1][Open --- needs the team]{%
  \begin{tcolorbox}[
    breakable,
    colback=white,
    colframe=black,
    colbacktitle=black!8,
    coltitle=black,
    boxrule=0.4pt,
    titlerule=0pt,
    left=0.6em,right=0.6em,top=0.5em,bottom=0.5em,
    title={#1},
    fonttitle=\bfseries\small,
    sharp corners,
  ]\small
}{%
  \end{tcolorbox}
}

\NewDocumentEnvironment{infobox}
  { O{Open --- needs the team} O{} }
{%
  \begin{tcolorbox}[
    breakable,
    colback=white,
    colframe=black,
    colbacktitle=black!8,
    coltitle=black,
    boxrule=0.4pt,
    titlerule=0pt,
    left=0.6em,
    right=0.6em,
    top=0.5em,
    bottom=0.5em,
    title={#1},
    fonttitle=\bfseries\small,
    sharp corners,
    #2
  ]\small
}{%
  \end{tcolorbox}
}


% Status markers for the submission checklist.
\newcommand{\stdone}{\ding{51}}
\newcommand{\stpartial}{\ding{108}}
\newcommand{\stgap}{\ding{55}}

% Title block.
\newcommand{\sostitle}[3]{%
  \begin{center}
    {\LARGE\bfseries #1}\\[0.35em]
    {\large\itshape #2}\\[0.5em]
    {\small #3}
  \end{center}
  \vspace{0.4em}
  {\color{sosRule}\hrule height 0.4pt}
  \vspace{0.9em}
}


\hypersetup{pdftitle={Project Journey Portfolio — StudyOS Thesis-Finder}}

\begin{document}

\sostitle{Project Journey Portfolio}{StudyOS Thesis-Finder}{%
  Group submission \quad\textbullet\quad Dominique, Max, Valentin \quad\textbullet\quad University of Tübingen}

\begin{tcolorbox}[colback=white,colframe=sosRule,boxrule=0.6pt,left=0.7em,right=0.7em,top=0.5em,bottom=0.5em]
\textbf{The journey in one paragraph.} We started with a matching platform for
thesis search, built as an app with a database behind it. Early conversations
with professors broke that premise, because the supply side will not feed a
central platform and a curated list rots the moment nobody is paid to keep it
alive. We reframed the problem from matching students to topics towards the
student's own cold start, moved to a database-free agent skill that searches the
live web and published it as an installable package that needs no server and no
account. The part that never happened is the test with students who are actually
looking. We ran a session in the course and then a public beta, and apart from
one written form from our supervisor nobody answered.
\end{tcolorbox}

% --------------------------------------------------------------------------- Section 1
\section{Initial Project Idea}

\begin{infobox}[We solve $X$ for $Y$ by doing $Z$]
We solve \textbf{topic discovery} for \textbf{thesis students} by \textbf{bringing open theses into one place}
\end{infobox}

\noindent Our initial idea was a full-stack web app for students and professors. The professors could easily list their recent thesis topics, while the students could 
search through the list of available topics and find their most suitable match. The motivation came from our own experience of searching and finding a suitable thesis topic. At the university, there are multiple different data sources like the official university 
website, Alma, and additionally most of the professors also have a private lab-based website. 
Our first hypothesis was that the major problem is the lack of a single source of truth. 
Therefore, we have planned and developed a first full-stack prototype, with a backend, a database, and a web frontend. Since we are in the age of AI, we also planned to integrate a fully custom AI agent with custom tool calls that could supervise students during the search process and the exploration of their own possibilities and goals for their thesis. 
There we fell into the trap that we, as software engineers, love to build complex systems, just because we can, not because it is actually needed.\\
Our initial target group was students and professors of the Department of Computer Science at the University of Tuebingen. 
We treated the students as the demand side and the professors as the supply side of our platform. After discussing the initial idea in the course, a lot of students thought that the idea was promising, because a lot of them faced the same problem in their bachelor studies. Common stories were that students were frustrated in the end because they didn't know about specific chairs, 
which might have been the perfect fit for their interests and skills. Another main point was that a lot of students reached out to professors, but didn't get a response.
At this point, we hadn't done the user research for the supply side yet, so we only had the students' perspective and the perspective of Professor Gehler.\\

To mitigate this discrepancy, we planned a user study for professors, where they could give us their point of view, without rating the actual idea itself. We wanted to understand the reasons behind the missing thesis list and whether it was intended or just fell off due to the lack of time.

While the whole application got stripped apart during the process of the course, the initial idea was a good starting point and a problem worth solving.



% --------------------------------------------------------------------------- Section 2
\section{Discovering and Understanding Users}
\label{sec:users}

\subsection*{What}
\begin{itemize}
  \item Two user groups: students
    are the demand side, professors/chairs are the supply side a platform would
    depend on. We are in the middle and wanted to satisfy both sides.
  \item We interviewed a substantial number of professors via E-Mail.  We sent 4 probe e-mails to professor we knew, then we automatically sent mails to all the professors in the department of computer science. This took place from late April till May.  We received a lot of feedback, which we documented and analyzed. The main insights are summarized below.
    The E-Mails can be found in the Appendix~\ref{app:prof-interview}. The professors were contacted via the official university e-mail addresses. We sent out two different variants of the e-mail, depending on whether the chair already had a public list esis topics or not. The professors were asked to give us their honest feedback on the current situation and the proposed idea.
   Main documented
    insights:
    \begin{itemize}
      \item Roughly half rejected a central platform outright and only about a
        third were conditionally open to it.
        Reason for rejection where that the professors think that fitting a student to a thesis topic is a personal process and that they want to get to know the student first.
        One important quote "the old system was the exact opposite: Listing theses was cumbersome. The list wasn’t maintained. And being among the few groups that listed on the central page usually just meant that the less engaged, less informed students would approach you first."
      \item One professor reacted negatively to our e-mail and thought they were too demanding. One important lesson is that when you send out a request to professors, you should be very careful with the wording and the tone of the message.
      \item A comparable topic-listing platform at the university had already
        failed for incentive reasons rather than UI reasons, because nobody wants
        to list and maintain open topics. This caused us to shift our product without surveying students, since the supply side is the bottleneck.
   
    \end{itemize}
    So we learned that the best way to please professors is to not ask them to do much at all. They can still go on and update their own website, but we don't need their input necessarily. 

  \item We also anticipated professors complaining about AI spam, so we instructed the skill to output a clear disclaimer in order to encourage students to atleast modify the generated e-mails and drafts.
 

  
\end{itemize}


\subsection*{So what}
\begin{itemize}
  \item The initial idea could not have told us that the \emph{supply side} is
    the bottleneck. Interviews revealed the platform premise fails regardless of
    how good the product is, because participation and freshness are incentive
    problems rather than engineering problems.
\end{itemize}

\subsection*{Now what}
\begin{itemize}
  \item  The main insight was that the we could not work with a central list, but should focus on what is publicly avaiblable and on the students's side.
  \item This brougt us to the idea of building an advisor that helps students, as detailed in the next section.
\end{itemize}

% --------------------------------------------------------------------------- Section 3
\section{Defining the Problem}


\subsection*{What}
We at first tried to find a problem that fit our rough goal a problem we had and we knew from our own student life. Finding the right bachelor or master thesis can be difficult. 
At first we thought that the problem was that there was no central list of topic. 
With the feedback of professors we knew that this was not  a solution to the problem of finding a suitable thesis topic for students and it would not have been maintainable either.

\subsection*{So what}
\begin{itemize}
  \item We therefore discussed mutliple pivots, however we still wanted to build an app,that would have been deployed. It would have been a full-stack app with a database and a backend.
  \item We discovered our new formulation of the problem : We solve the problem of finding a suitable thesis topic for students by providing a tool that helps students find and reflect on their thesis by proving information and idea via an LLM.
  \item For example the app would be able to ingest the transcript of records and then try to build a profile that would be used to find suitable thesis topics.
\end{itemize}

\subsection*{Now what}
\begin{itemize}
  \item  We now had a clear idea of the problem and started building an app. We wanted to asses the students based on different axis, like interest or skill. We also scraped papers by the professors and their chairs. This was all done so we might use local models or would have a big database that we could use to answer the studetents questions.
  \item You can find the app in the git logs of our repository. However since we later decided to not use an app but a skill and then even later pivoted to a skill without any database ( even in text form) whatsoever, we have deleted it from the current iteration of the repository.
  \item What stuck from this phase was the main idea : Advising students. Also we carried the idea of building a profile of the student first
\end{itemize}



% --------------------------------------------------------------------------- Section 4
\section{Proposing a Solution}

\noindent As described in Section~3, we have changed the problem definition from a tool for professors and students to a tool for students only. After we've evaluated the answers from the professors, it became quite clear that they are not interested in the centralized app solution. Mainly because they don't want the overhead, and they want to get to know the student to find a suitable thesis.
Since the initial idea was built on the supply of the professors, we had to rethink our approach, because the problem for the students still exists. The core problem is not the absence of a central thesis-topic list. Students often lack the context to decide which chair or company fits their interests, methods, skills, constraints, and thesis style. This makes thesis search highly individual: students need guided matching, not just another static list. Since this problem is highly individual and needs deep information about the user, we thought about a private assistant.


\noindent We initially stuck to the idea of a full-stack app, with a parser that extracts the information from the professors' websites and then presents it to the user. After we presented our idea, the resulting reaction was again: "Why not just a Skill?" For us, it was quite new that we "just" use a markdown file to describe the task and an LLM will do the rest. 
After further discussion in the team, we really embraced the idea of a Skill, because a Skill can be inspected and adapted through plain Markdown, lowering the barrier for future maintainers.


Additionally, we don't have to think about security, privacy, deployment, and cloud billing since this is already handled by your LLM provider of choice. This realization then led to the initial 
implementation of our first skill architecture, with a markdown database of all possible professors, including their recent papers, teams, and their research topics. The skill then used this information to answer student-specific questions in a conversational way. 
During this phase, we thought deeply about what information is actually necessary to provide high-quality answers. We have developed a 6-dimensional scheme, which needs to be answered before the student gets any recommendation from the LLM. The 6 dimensions are

\begin{table}[H]
  \centering
  \begin{tabular}{|c|l|p{10cm}|}
    \hline
 Dimension & Name & Description \\
    \hline
 1 & Interests & Core research areas or topics \\
    \hline
 2 & Methods & How the student wants to work, e.g. empirical, qualitative, computational \\
    \hline
 3 & Domain & Application field, e.g. healthcare, education, finance, automotive \\
    \hline
 4 & Thesis style & Preferred output type, e.g. experimental, theoretical, systems, analysis, survey, mixed \\
    \hline
 5 & Skills & Concrete tools or competencies, e.g. Python, ML frameworks, lab methods \\
    \hline
 6 & No-gos & Hard exclusions, e.g. hardware setup, pure proofs, large-scale software engineering \\
    \hline
  \end{tabular}
\end{table}

\noindent In our opinion, those 6 dimensions are the most important cornerstones for a successful thesis search.

\noindent The biggest task here was "Why should a student use our skill and not just plain Claude/Codex?" The answer to 
this question is that our skill is not a "give me the answer as fast as possible", but it emulates a session with a professional study coordinator that explicitly tries
to find the best possible match for the individual student.


\noindent After we have implemented this orchestrated multi-dimensional skill, we've tested it thoroughly and recognized that the markdown database needs to be maintained in some way, which we wanted to avoid in the first place. We extended our current prototype and built a database-less approach, where we are utilizing
the commonly used web search tool of the LLM providers to find the most suitable professors and their research topics. This approach has the advantage that we don't have to maintain a database, which can be costly and error-prone if, e.g., new Professors join or old ones leave the Department of Computer Science. 
The biggest disadvantage here was that the search tool introduces even more uncertainty and non-deterministic behavior into the skill. To mitigate this, we have built explicit search templates to get the most relevant information from the web search.
We conducted a field study and let the course test
both skills, the one with the markdown database and the one without. The feedback from that session was spoken rather than written, and it pointed at the variant without the database as the more useful one. This led to the final decision of removing the Markdown database and using the slimmer version with the minimal maintenance overload. The session itself is described in Section~6.
Another finding during the field study was that some students also want to do a thesis in a company. Therefore, we extended the skill to also provide suitable thesis topics from companies in Baden-Württemberg.



\noindent Since the skill is more uncertain, and a session takes longer than 30 minutes, we decided to use some automated testing. Therefore, we built evals as well as an orchestrated skill set to run end-to-end simulations with an LLM-as-a-judge to evaluate the actual performance. These automated tests gave a major boost in confidence and implementation speed,
since the agent can now fully simulate actual conversations.
In order to make this representative, we designed multiple fictitious students
from different backgrounds and with different skill sets and answer styles. 


\noindent Our final product is a low-maintenance, database-free skill package that guides students from vague thesis interests to evidence-backed university or company options, without requiring professors to maintain yet another platform.
% --------------------------------------------------------------------------- Section 5
\section{Building}

\noindent The product itself is submitted separately in the Product Artifact ZIP, so this section only refers to that submission.

\noindent The build followed the priority order that came out of Section~4. We started with the university arm, added the company arm on the same database-free principle and consolidated both into a single entry-point skill, so that a student does not have to decide up front which of the two paths they need.

% --------------------------------------------------------------------------- Section 6
\section{Testing}

\subsection*{What}

\noindent We tested the skill ourselves throughout, using the invented student
profiles from Section~4. The profiles from medicine and neurosciences were the
most useful, because we have no background there and a weak answer shows
immediately.

\noindent We also compared the skill against chair lists we had curated
ourselves and against a plain assistant. We do not report those numbers, because
we wrote both the ground truth and the scoring, and our first design graded the
skill on material it had already seen. We fixed that, but the benchmark is still
our own, so we use these runs as a regression check only.

\noindent A handful of friends and fellow students ran the skill in front of
us, most of them from computer science. Nothing was recorded, so this is
reported from memory. Reactions were
positive and several said they would treat the suggestions as a starting point
rather than act on them. None of them were choosing a thesis at the time.

\noindent The one test outside that circle was the course session in early July,
where each team presented its prototype and let the room try it. We brought both
variants, with and without the markdown database. Spoken feedback was positive
and the point raised most often was that the tool needs no installation and no
account, since it runs inside an assistant people already use. The same
conversations favoured the database-free variant, which is the feedback
Section~4 refers to. Only the course professor answered in writing.

\noindent He found the results useful, the most useful one being that he had
``found a new colleague I did not know about''. The interaction he found harder.
The session ``was asking a lot'' without making clear how much detail was
expected at which step, and it stopped where he would have wanted help going
deeper into a single option. The
form is in Appendix~\ref{app:gehler-feedback}.

\noindent We then published the skill package with an install guide, sent a
recruiting mail in German and English and linked an anonymous survey. The survey
asked whether the run had changed how clearly the student sees their own thesis,
plus relevance, source trust and whether they would contact anyone suggested. We
asked people whose run had broken off to answer anyway. We received no
responses. The mail and the survey are in
Appendix~\ref{app:beta-mail}.

\noindent We also asked Professor Hennig, who had answered our first interview
round with the most enthusiasm, what he made of the finished tool.
He has not replied so far.

\subsection*{So what}

\noindent The feedback confirmed the two decisions the project rests on. The
live search finds things a real user did not know, which is what the
database-free approach was meant to deliver, and the property people valued most
in the course was that nothing has to be set up, which is what the skill format
was chosen for. What the feedback added is where the remaining work is. The
interview needs to say why it asks what it asks, and starting a session has to
become as easy as people assumed it already was.

\subsection*{Now what}

\noindent Three changes shipped within a day of the course session. The skill
states up front what it will ask and why. It offers to pick one option and go
deeper instead of stopping at the option map. Drafting a contact mail now
requires reading at least one paper of the chair first. Two requests were
dropped. Letting professors supply input would have rebuilt the supply-side
dependency the project had just left, and an example walkthrough would not
have fixed the session ending too early, which was the actual problem.

% --------------------------------------------------------------------------- Section 7
\section{Iteration and Next Steps}

\subsection*{What}

\noindent Two changes from feedback changed what the product is. The
first was the pivot described in Sections~2 and~3, when the professor interviews
ended the platform premise and with it the full-stack application we had already
built. The second came in August, when we asked our supervisor what mattered more
to him, the proposals the skill generates or the fact that a student had used it
to think about what they want to do. He answered that the reflection mattered
more. The skill now records the student's own description of the thesis they are
looking for before the search and again afterwards and shows both side by side.
The three changes from the course test are described in Section~6. What has
not changed is that the company arm is weaker than the university arm, since
companies in Baden-W\"urttemberg publish far less about what they would
supervise than chairs do.

\noindent The skill package is published as a versioned release on GitHub with
an install guide, and we tested a cold install from the archive alone. It is
meant to go into a repository of the AI Center, which would give it an address
that does not depend on us, and we have asked the student advisory service
whether it could be mentioned at the information events on bachelor theses,
where students with exactly this question are. That mail is in
Appendix~\ref{app:seibold-mail} and we are waiting for an answer.

\noindent Maintenance is a question we answered by design. The product is a
folder of Markdown files with no database, no backend, no server, no credentials
and no build step, so the failure modes that usually end a student project do
not apply. The repository lives in the \texttt{Tue-StudyOS} organisation on
GitHub rather than in a private account. Three things can still go stale, which
are the skill format if the providers change it, the university websites if they
are restructured and our discovery rules if students start searching
differently. None of these needs somebody to watch the project, only somebody to
pick it up on the day it stops working, which we estimate at a couple of hours
per semester. In that sense the skill is finished, and that is why we do not
name an owner.

\noindent This is also where AI has the clearest role. The repository ships an
operating guide written for future AI agents rather than for people, in
\repofile{AGENTS.md}, so the intended maintainer of a product that was largely
built with agents is itself an agent, with a human spending an hour or two per
semester steering it and checking the result.

\noindent The last thing we did was test whether any of this is specific to
theses. The skill is a sequence that does not depend on the subject. It
interviews the person until a profile exists, searches live sources against that
profile and hands back options with evidence. So we rebuilt it for career
advice, at \url{https://github.com/ValentinJSchmidt/german-career-compass-skills}.
It took a fraction of the original effort, because a coding agent can adapt an
existing structure from a few examples, and it worked better than we expected.
It is an early version, but it shows that the transferable result of this
project is the sequence and not only the thesis finder built on top of it.

\subsection*{So what}

\noindent The improvements came from two sources. Outside input showed us the
problem. The professor interviews ended the platform, and a short presentation
in the course ended the web app, because the reaction was a question we had not
asked ourselves, which was why we were not simply using a skill. In both cases
the value was the different vantage point. Our own thinking produced the
constraints. Dropping the markdown database came from asking who would keep it
up to date and concluding that the capacity for that would not exist. Nobody
outside the team raised it. Outside input told us which problem to solve, and
asking who maintains this in three years told us how to build it.

\subsection*{Now what}

\noindent The next steps are in order. First lower the cost of starting a
session, since Section~6 shows that this is the barrier. Then run the tool with
students who are currently looking for a thesis, which is the test we have
prepared but not yet done. The product itself is complete. What it needs is a reason for somebody to open it,
and that is a distribution question rather than an engineering one.

% --------------------------------------------------------------------------- Appendix
\appendix
\section{Interview Guide --- Professors}
\label{app:prof-interview}

\noindent The professor research referenced in Section~\ref{sec:users} was
conducted by e-mail. We sent two variants depending on whether the chair already
published thesis topics or not; both are reproduced verbatim below.


\noindent\textbf{Mode:} e-mail \quad\textbullet\quad
\textbf{Period:} \emph{[dates]} \quad\textbullet\quad
\textbf{Contacted:} \emph{[n]} chairs \quad\textbullet\quad
\textbf{Responses:} \emph{[n]}

\subsection*{Mail A --- chair already lists topics publicly}

\noindent\textbf{Subject:} Thesis topic listings at your chair -- your experience \& a new approach

\medskip
\noindent Dear [Name],

\medskip
\noindent We are Master's students in Machine Learning at the University of
T\"ubingen, working on a project supervised by Prof.\ Peter Gehler to improve how
thesis topics are shared and managed across chairs in our department.

\medskip
\noindent We noticed that your chair already lists thesis topics publicly ---
which is genuinely rare and makes your perspective especially valuable to us.

\medskip
\noindent We have a few concrete questions about your current experience:

\begin{enumerate}[leftmargin=1.4em,itemsep=0.35em]
  \item How much effort does it take to keep your topic list up to date --- is
    it manageable, or does it create noticeable overhead?
  \item Do topics sometimes stay listed past their actual availability, or get
    delayed because updating feels like extra work?
  \item How do you currently decide when and how to update the list --- is there
    a process, or does it happen ad hoc?
\end{enumerate}

\noindent We're developing a concept we'd love your honest reaction to: a
centralized platform where chairs manage their topics and receive periodic
reminders to review and update their listings --- keeping the overhead as low as
possible.

\medskip
\noindent A few follow-up questions on that:

\begin{enumerate}[leftmargin=1.4em,itemsep=0.35em,resume]
  \item Is that something your chair would actually want to use --- and if not,
    what would need to be different for it to make sense?
  \item Are there concerns around data privacy or how a shared system would
    handle your content that we should be aware of?
\end{enumerate}

\noindent A brief written reply would already help us enormously. If anything is
easier to discuss directly, we're happy to find a time.

\medskip
\noindent Best regards, [Your name] (Project supervised by Prof.\ Peter Gehler)

\subsection*{Mail B --- chair has no public topic listings}

\noindent\textbf{Subject:} Thesis topics at your chair -- a few questions about going public

\medskip
\noindent Dear [Name],

\medskip
\noindent we are Master's students in Machine Learning at the University of
T\"ubingen, working on a project supervised by Prof.\ Peter Gehler on how thesis
topics are currently --- and often not --- shared across chairs in our
department.

\medskip
\noindent We noticed that your chair doesn't currently list thesis topics
publicly, and we're genuinely trying to understand why before drawing any
conclusions. All answers are useful to us, including ``we don't want to.''

\medskip
\noindent A few concrete questions:

\begin{enumerate}[leftmargin=1.4em,itemsep=0.35em]
  \item Is the absence of public listings a deliberate choice --- for example,
    because you prefer students to reach out directly --- or more a matter of no
    good tool existing to make it easy?
  \item Are there topics internally that would in principle be suitable for
    students, or do topics at your chair tend to emerge from ongoing research
    rather than being defined in advance?
  \item If maintaining a public list has felt like too much overhead in the
    past: what specifically makes it costly --- writing the descriptions,
    keeping them current, or something else?
\end{enumerate}

\noindent We're developing a concept we'd love your honest reaction to: a
centralized platform where chairs manage their topics and receive periodic
reminders to review and update their listings --- keeping the overhead as low as
possible.

\begin{enumerate}[leftmargin=1.4em,itemsep=0.35em,resume]
  \item Is that something your chair would actually want to use --- and if not,
    what would need to be different for it to make sense?
  \item Are there concerns around data privacy or how a shared system would
    handle your content that we should be aware of?
\end{enumerate}

\noindent We're not trying to push a solution --- we want to know whether a
problem worth solving actually exists here from your side. A brief written reply
is everything we need. Happy to chat directly if that's easier.

\medskip
\noindent Best regards, [Your name] (Project supervised by Prof.\ Peter Gehler)


% --------------------------------------------------------------------------- Appendix B
\section{Course Test --- Written Feedback}
\label{app:gehler-feedback}

\noindent Verbatim answers from the feedback form after the in-course prototype
session in early July. This is the only completed form we received.

\medskip
\noindent \textbf{How natural did the interaction feel?}\\
``It was asking a lot. It was not clear how much detail was required at which
step''

\medskip
\noindent \textbf{Was it clear why certain chairs, papers, companies, or thesis
directions were suggested?}\\
``Yes, very good. Was unexpectedly good. This was good''

\medskip
\noindent \textbf{How trustworthy did the suggestions feel?}\\
``Good!''

\medskip
\noindent \textbf{How well did the skill understand what you are looking for?}
(1--5)\\
3

\medskip
\noindent \textbf{How personally relevant did the recommendations feel?}
(1--5)\\
4

\medskip
\noindent \textbf{What was the most useful output of the skill?}\\
``Found a new colleague I did not know about.''

\medskip
\noindent \textbf{What was missing for you to take a real next step after the
session?}\\
``I think a bit more follow up help to be offered. Like tell me which one you
would choose and then walk a bit more down that road. Eg also a bit more clarity
of what I likely learn.''

\medskip
\noindent \textbf{Would you use the skill again or recommend it to another
student?}\\
``Can a professor provide information for the skill. Eg looking for \dots\ Not
looking for\dots\ / An example walkthrough, or like \,``Lisa found her thesis
\dots'' / Example prompts? Maybe explain what info is needed or give
suggestions''

% --------------------------------------------------------------------------- Appendix C
\section{Beta Test Recruiting Mail}
\label{app:beta-mail}

\noindent Sent in German and English. The English version is reproduced here. It
produced no responses.

\medskip
\noindent Hi everyone,

\medskip
\noindent We built a set of AI ``skills'' that help you find a thesis topic and
supervisor, and we need people to actually try them out and tell us what
happened.

\medskip
\noindent \textbf{What it does.} You type one command and it interviews you about
your interests, methods and constraints, then searches the live web for
University of T\"ubingen chairs or Baden-W\"urttemberg companies that fit, and
hands you a map of real options with contact paths. The point is less the list
than the interview. Most people come out with a sharper idea of what kind of
thesis fits them than they went in with.

\medskip
\noindent \textbf{What we need from you:}

\begin{enumerate}[leftmargin=1.4em,itemsep=0.35em]
  \item Install it. Instructions for every route (Claude, ChatGPT, Gemini, or a
    terminal agent) are at \url{https://github.com/Tue-StudyOS/study-os-thesis/blob/main/INSTALL.md},
    downloads at \url{https://github.com/Tue-StudyOS/study-os-thesis/releases/latest}.
    No account, no signup, nothing installed outside one folder.
  \item Run it. Type \texttt{thesis-finder} and answer the questions. Plan for
    20--40 minutes. You need web search enabled in your assistant.
  \item Fill in the survey, about 4 minutes:
    \url{https://www.soscisurvey.de/thesis-skill/}
\end{enumerate}

\noindent We do not receive any of your conversations.

\medskip
\noindent Please answer the survey even if the run went badly or you broke it off
halfway. A failed run is a result for us, and those are the answers that actually
change the tool.

\medskip
\noindent The survey is anonymous, we don't store IP addresses, and you can stop
at any point.

\medskip
\noindent Thanks a lot. Every single response helps.

\medskip
\noindent Best, Valentin, Maximillian and Dominique

\subsection*{The linked survey}

\noindent Anonymous, hosted on SoSci Survey on German servers with a consent
page and no IP logging, about four minutes. It received no responses. The
consent text is omitted here. All batteries are five-point agreement scales.

\medskip
\noindent \textbf{Before the session compared with now}

\begin{itemize}[leftmargin=1.4em,itemsep=0.2em]
  \item My idea of my thesis is just as vague now as it was before.
  \item I have a clearer idea of what kind of thesis would fit me.
  \item Since the session I have thought about my thesis more concretely than before.
  \item The session made me reconsider something I thought I wanted.
  \item Everything it showed me I already knew about.
  \item I came across directions or research groups I would not have found on my own.
\end{itemize}

\noindent \textbf{Duration} of the whole run, from the first question to the
final list. Less than 15 minutes, 15 to 30, 30 to 45, 45 to 60, more than an
hour.

\medskip
\noindent \textbf{The chairs, groups or companies shown} (with an additional
\emph{I cannot judge this} option)

\begin{itemize}[leftmargin=1.4em,itemsep=0.2em]
  \item The suggestions were concrete enough to act on.
  \item I intend to contact at least one of the people or places suggested.
  \item I would trust the information enough to write to one of them.
  \item Please select ``disagree'' here so we can tell careful answers apart from
    clicked-through ones. (attention check)
  \item The suggestions fit what I am actually interested in.
  \item Most of the suggestions felt generic and could have gone to any student.
\end{itemize}

\noindent \textbf{The process itself}

\begin{itemize}[leftmargin=1.4em,itemsep=0.2em]
  \item Answering all the questions felt tedious.
  \item The questions I was asked were relevant to me.
  \item It felt like it understood what I actually care about, not just my keywords.
  \item I would recommend it to a friend who is looking for a thesis topic.
  \item The time it took was worth what I got out of it.
\end{itemize}

\noindent \textbf{Free text.} What was the single most useful thing about it?
What annoyed you, or what did it get wrong? Respondents whose run stopped
partway were asked to answer anyway and say where it broke off.

% --------------------------------------------------------------------------- Appendix D
\section{Mail to the Student Advisory Service}
\label{app:seibold-mail}

\noindent Sent in German and translated here. No answer so far.

\medskip
\noindent Dear Ms Seibold,

\medskip
\noindent you will already be familiar with the StudyOS course, and like other
groups we are now approaching you. We have built several skills for AI agents
that help a student find a thesis, in T\"ubingen.

\medskip
\noindent We hope that Prof.\ Gehler will take these skills into a repository of
the AI Center. It would mean all the more to us if students actually used the
work, because we believe the skills can be a useful building block in finding a
suitable thesis.

\medskip
\noindent We have also taken care not to make it too easy for students, so that
reflection actually happens and teaching staff are not flooded with AI slop.

\medskip
\noindent Would it be possible to mention these skills at the information events
on theses? Or would we have to contact somebody else for that?

\medskip
\noindent Thank you for communicating so openly with students.

\medskip
\noindent Best regards, Valentin Schmidt, Maximilian Schnitt, Dominique Meyer


\end{document}
